\section{Smooth ambient germs and their ordinary orders}
\label{sec:ambient-germs}

Starting from the two graph maps, we now construct the smooth fourfold
germs and compute the first nonzero terms of their functions.  We first
fix one parameter choice for both graphs.  The same choice will satisfy
the later arithmetic requirements, while the estimates proved here
supply the local isolation argument in Lemma~\ref{lem:isolation}.

\begin{lemma}[One common parameter choice]\label{lem:simultaneous-parameters}
For every sufficiently large integer \(r\), choose an integer \(n_0\)
such that
\[
 n_0\equiv1\pmod{2^r},\qquad n_0\equiv2\pmod5,\qquad
 n_0>5\cdot2^r,
\]
and define
\begin{align*}
 n_1&=n_0+5\cdot2^r,&
 a&=(n_0^2-1)/4,& b&=(n_1^2-1)/4,\\
 \alpha&=\frac{4b}{a+b},&
 \beta&=\frac{4a}{a+b}.
\end{align*}
Then \(a,b\) are positive comparable integers,
\(n_0,n_1\equiv1\pmod4\), and \(\gcd(n_0,n_1)=1\).
Moreover
\[
 \alpha+\beta=4,\qquad a\alpha=b\beta,\qquad
 \frac{(a-2)\alpha}{4}>4,\qquad
 \frac{(b-2)\beta}{4}>4.
\]
\end{lemma}

\begin{proof}
The Chinese remainder theorem supplies arbitrarily large \(n_0\)
with the two congruences, so the strict lower bound can also be imposed.
For large \(r\) we have \(r\ge2\), and both \(n_j\) are \(1\pmod4\);
therefore \(a,b\) are integers.  The lower bound on \(n_0\) gives
\(n_0<n_1<2n_0\), so \(a,b\) are positive, comparable, and grow
without bound with \(r\).  A common divisor of \(n_0,n_1\) divides
\(5\cdot2^r\).  Both integers are odd and \(2\pmod5\), proving
\(\gcd(n_0,n_1)=1\).

The identities for \(\alpha,\beta\) follow from their definitions.
Since \(b>a\),
\[
 \frac{(a-2)\alpha}{4}
 =\frac{ab}{a+b}-\frac{2b}{a+b}
 >\frac a2-2,\qquad
 \frac{(b-2)\beta}{4}
 =\frac{ab}{a+b}-\frac{2a}{a+b}
 >\frac a2-2.
\]
The right sides exceed four for sufficiently large \(r\).
In particular, the chosen \(a,b\) exceed four.
\end{proof}

Fix this choice for the rest of the paper, and set
\[
 g=\gcd(a,b),\qquad m=4ab+a+b.
\]
Lemma~\ref{lem:parameter-indices} proves the additional arithmetic
properties of these same integers: \(a,b\equiv2\pmod5\),
\(5\nmid g\), \(g\ge2^{r-1}>5\), and, for each \(q\in\{1,4\}\),
the existence of positive integers \(i,j\) satisfying
\[
 i\equiv j\pmod4,\qquad \frac{bi+aj}{m}=\frac q5.
\]
The deformation parameter \(s\) will be chosen sufficiently small only
after this tuple has been fixed.

Put \(p=(X,T)\).  Let \(R_0(p,w)\) and \(R_1(p,w)\) be the monic
degree-four Weierstrass polynomials in the graph equations
\begin{align*}
 X-T^2-q_1(T,w)&=U_0(p,w)R_0(p,w),\\
 T-\psi(X,w)&=U_1(p,w)R_1(p,w),
\end{align*}
where \(U_0,U_1\) are analytic units.  For a fixed small real \(s>0\),
set
\begin{align*}
 R_{i,s}(p,w)&=s^4R_i(p,w/s),\\
 C_{i,s}&=\{(p,y,z,w):yz=R_{i,s}(p,w)\},\\
 f_{i,s}&=\left.h(p)+w(y^a+z^b)\right|_{C_{i,s}} .
\end{align*}

\begin{proposition}[Analytic germs and their orders]
\label{prop:ambient-orders}
For the parameters of Lemma~\ref{lem:simultaneous-parameters}
and every sufficiently small \(s>0\), each \(C_{i,s}\) is a smooth
complex fourfold near the origin, and \(f_{i,s}\) has an isolated
critical point there, as proved in Lemma~\ref{lem:isolation}. After solving the smooth ambient equations
in analytic local coordinates, \(f_{0,s}\) has order five and
\(f_{1,s}\) has order four. Both are reduced convergent germs.
\end{proposition}

\begin{proof}
The specializations in Lemma~\ref{lem:plane-data} give
\(R_i(0,w)=w^4\) and \(U_i(0)=1\).  Hence
\[
 \partial_XR_0(0)=1,\qquad \partial_TR_1(0)=1.
\]
These derivatives remain nonzero along the respective local
root loci after shrinking.  For \(s>0\), the corresponding derivatives
of \(R_{i,s}\) are multiplied by the nonzero factor \(s^4\).
At a point of \(C_{i,s}\) with
\(y\ne0\) or \(z\ne0\), its defining equation has a nonzero
vertical derivative; at \(y=z=0\), one of the displayed base
derivatives is nonzero. Thus the ambient germs are smooth.

The two ambient equations are equivalently the implicit equations
\begin{align*}
 X&=T^2+q_1(T,w/s)+s^{-4}U_0(p,w/s)yz,\\
 T&=\psi(X,w/s)+s^{-4}U_1(p,w/s)yz,
\end{align*}
respectively.  In the first, \(q_1(T,w/s)\) has ordinary order at least
three.  Since \(U_0(0)=1\), analytic implicit solution gives
\[
 X=T^2+s^{-4}yz+O_{\ge3}(T,y,z,w).
\]
In the second graph \(q_2=q_1-u^2\) has order two, and
\(\psi=u(u-w^2)q_2\) has order four. Therefore
\[
 T=s^{-4}yz+O_{\ge3}(X,y,z,w).
\]
Here \(s\) is fixed and \(O_{\ge j}\) denotes terms of ordinary
total degree at least \(j\); in particular all \(w^4\) terms
lie in these remainders. The chosen \(a,b\) exceed four, so
\(w(y^a+z^b)\) has order greater than five. Substitution in \(h\)
gives the nonzero initial forms
\[
 \operatorname{in}(f_{0,s})
   =T(s^{-4}yz)(s^{-4}yz-T^2),\qquad
 \operatorname{in}(f_{1,s})=s^{-4}yzX^2.
\]
They have degrees five and four, respectively.

It remains to justify the isolated critical points claimed in the
proposition.  The input from the base curve is the following explicit
gradient estimate.  Along an analytic arc, let \(\nu\) denote order in
the arc parameter, with \(\nu 0=\infty\); for a vector, \(\nu\) is the
minimum of its component orders.  If an analytic arc \(p=(X,T)\ne0\) has
\(m_0=\nu p=\min(\nu X,\nu T)\), then
\begin{equation}\label{eq:plane-gradient}
 \nu\nabla h\le4m_0.
\end{equation}
Indeed
\(h_X=T(2X-3T^2)\) and
\(h_T=X^2-9T^2X+10T^4\).
If \(\nu X<2\nu T\), the first term of \(h_T\) has order
\(2\nu X\le4m_0\). If \(\nu X>2\nu T\), the term
\(-3T^3\) of \(h_X\) has order \(3\nu T\le4m_0\).
In the equality case, put \(X/T^2=c+\text{higher terms}\);
\(2c-3\) and \(c^2-9c+10\) have no common root, so one
derivative has order at most \(4\nu T=4m_0\).
The cases in which one coordinate vanishes identically are
immediate.

The full multiplier argument, including the cases where a vertical
variable vanishes, is given in Lemma~\ref{lem:isolation}; it uses the
two strict bounds in Lemma~\ref{lem:simultaneous-parameters} with the
gradient exponent \(4\) just proved.

Finally, if a convergent hypersurface germ had a repeated
irreducible factor, its gradient would vanish at generic
nonzero smooth points of that factor. Lemma~\ref{lem:isolation} therefore also implies reducedness.
\end{proof}

Choose one \(s>0\) below the thresholds for both graphs and choose
analytic charts \(C_{i,s}\cong(\mathbb C^4,0)\).  The germs in
Theorem~\ref{thm:main} are then \(f_5=f_{0,s}\) and \(f_4=f_{1,s}\).
